\subsection{11.05.2026 (Freistunde)}

\subsubsection{Ziel}

Ziel des Tages war es, eine Webseite zur
Darstellung und Steuerung der Taupunktlüftung
zu erstellen.

\subsubsection{Durchgeführte Arbeiten}

Zunächst wurde eine einfache Webseite mit dem
Python-Framework Flask erstellt.

Dabei wurde eine Startseite implementiert,
welche eine kurze Erklärung zur Funktion der
Taupunktlüftung und der Webseite enthält.

Zusätzlich wurde eine weitere Seite erstellt,
auf der die letzten 50 gespeicherten Datensätze
aus der Datenbank tabellarisch dargestellt werden.

Die Datenbanktabelle wurde dafür um den
Lüfterstatus erweitert, damit sichtbar ist,
ob der Lüfter zum jeweiligen Zeitpunkt aktiviert
oder deaktiviert war.

Darüber hinaus wurde eine zusätzliche Seite
implementiert, auf der sich ein Button befindet,
mit dem der Lüfter manuell ausgeschaltet werden
kann.

Die manuelle Deaktivierung gilt dabei bis zur
nächsten automatischen Berechnung der
Taupunktlogik.

Während der Entwicklung auftretende Fehler
wurden analysiert und mithilfe von
:contentReference[oaicite:0]{index=0}
behoben.

\subsubsection{Erkenntnisse}

Es wurde festgestellt, dass Flask eine einfache
Möglichkeit bietet, Python-Anwendungen mit einer
Weboberfläche zu erweitern.

Zusätzlich wurde deutlich, dass sich Daten aus
SQLite problemlos in Webseiten integrieren und
tabellarisch darstellen lassen.

Außerdem wurde erkannt, dass eine Weboberfläche
die Bedienung und Überwachung der
Taupunktlüftung deutlich vereinfacht.

\subsubsection{Nächste Schritte}

Als nächster Schritt soll die Webseite getestet
und gegebenenfalls weiter angepasst werden.